\section{Unconditional theta quasimodes and the bottom eigenline}
\label{sec:ACspectralboundary}

The next results refine the dependence on the input theorem \(A\).
They use the unconditional closed, lower semibounded finite-window
operators.  In particular, a large subspace at energies very close to
zero is available without knowing the sign of the bottom eigenvalue.
Throughout, \(\widehat f(t)=\int f(x)e^{-itx}\,dx\), \(L=\log\lambda\),
and \(I_L=[-L,L]\).  The normalized theta kernel \(K\) is AC-S04.

\subsection{The global distribution identity}

Put
\[
 a(t)=-\log\pi+
       \operatorname{Re}\psi(1/4+it/2),
 \qquad \psi=\Gamma'/\Gamma .
\]
The centered Weil distribution in logarithmic coordinates is
\[
\begin{split}
 \mathcal W(F)={}&
 \frac1{2\pi}\int_{\mathbb R}a(t)\widehat F(t)\,dt
 -\sum_{n\geq2}\frac{\Lambda(n)}{\sqrt n}
                   \bigl(F(\log n)+F(-\log n)\bigr)\\
 &+\int_{\mathbb R}F(x)(e^{x/2}+e^{-x/2})\,dx .
\end{split}                                                    \tag{SB.1}
\]
Here \(\Lambda(n)\) is the von Mangoldt function.  The exact source
locators are S21, equations \texttt{bombtest},
\texttt{thetaprime}, \texttt{bombieriexplicit1bis}, and
\texttt{bombieriexplicit2bis}; see also (T.11) for its unconditional
zero-side identity.  The polar terms in (SB.1) are retained.

\begin{lemma}[AC-S22: the theta kernel is an unconditional null function]
\label{AC-S22}
Let \(\mathcal T\) be the space of smooth functions for which
\[
 \sup_x e^{b|x|}|F^{(j)}(x)|<\infty
 \quad\text{for every }b>0,\ j\geq0.
\]
All three terms in (SB.1) define a continuous functional on this
space.  The unconditional explicit formula extends to \(\mathcal T\).
In particular, for every integer \(j\geq0\),
\[
 \mathcal W*K^{(j)}=0, \qquad
 (\mathcal W*u)(x):=\mathcal W_y\bigl(u(x-y)\bigr).
                                                               \tag{SB.2}
\]
This is a distribution identity with a well-defined smooth input,
not the assertion that a whole-line closed Weil operator has been
constructed on unweighted \(L^2\).
\end{lemma}
\begin{proof}
The archimedean multiplier satisfies
\(|a(t)|\leq C\log(2+|t|)\); thus its term is continuous even on
the Schwartz space.  The polar term is continuous in a weighted
\(L^1\) norm.  For \(b>1/2\), the prime term is absolutely bounded by
\[
 2\sup_x e^{b|x|}|F(x)|
       \sum_{n\geq2}(\log n)n^{-b-1/2}<\infty.
\]
This proves the first assertion.

For a nontrivial zero \(\rho\), put \(z_\rho=i(\rho-1/2)\).
The explicit formula on compact smooth tests is
\(\mathcal W(F)=\sum_\rho m_\rho\widehat F(z_\rho)\).
Every \(z_\rho\) has imaginary part in \([-1/2,1/2]\).
Two integrations by parts give
\[
 |\widehat F(z)|
 \leq \frac{C}{1+|z|^2}
       \sum_{k=0}^2\int_{\mathbb R}
                       e^{|x|/2}|F^{(k)}(x)|\,dx
 \quad(|\operatorname{Im}z|\leq1/2).
\]
The unconditional Riemann--von Mangoldt bound implies
\(\sum_\rho m_\rho(1+|z_\rho|^2)^{-1}<\infty\).
The zero side is therefore continuous on \(\mathcal T\).
Smooth compact cutoffs tend to any \(F\in\mathcal T\) in all these
seminorms, proving the extended identity.

For \(F_x(y)=K^{(j)}(x-y)\), AC-S04 gives
\[
 \widehat F_x(z)=e^{-izx}(-iz)^j\Xi(-z).
\]
At \(z=z_\rho\), the last factor is
\(\xi(1/2-iz_\rho)=\xi(\rho)=0\).
Every summand vanishes, with no assumption on the real part of
\(\rho\).  This proves (SB.2).
\end{proof}

\subsection{The sharp cutoff belongs to the operator domain}

\begin{theorem}[AC-S23: an unconditional operator residual estimate]
\label{AC-S23}
For every fixed \(j\geq0\), put
\[
 f_{j,L}=1_{I_L}K^{(j)},\qquad
 r_{j,L}=(1-1_{I_L})K^{(j)},\qquad d_j=9/2+2j.
\]
Then \(f_{j,L}\in\operatorname{Dom}A_L\), and, as
\(\lambda=e^L\to\infty\),
\[
 \|A_Lf_{j,L}\|_2
     =O_j\!\left(\lambda^{d_j}(\log\lambda)^{3/2}
                              e^{-\pi\lambda^2}\right).         \tag{SB.3}
\]
The estimate is unconditional and includes both endpoints, the
infinite prime tail, and the two polar terms.
\end{theorem}
\begin{proof}
Let \(a(D)\) be the self-adjoint Fourier multiplier \(a(t)\) on the
whole line.  A function \(v\) in \(L^1\) with finite total variation
satisfies
\[
 |\widehat v(t)|\leq
       \min(\|v\|_1,\operatorname{Var}(v)/|t|).
                                                               \tag{SB.4}
\]
Consequently it belongs to \(\operatorname{Dom}a(D)\) whenever it
also belongs to \(L^2\).  Indeed
\(\int\log^2(2+|t|)\min(1,|t|^{-2})\,dt<\infty\).
The sharp cutoff \(f_{j,L}\) has these properties.  The finite-window
form is the restriction of the \(a(D)\) form plus the bounded
translation and polar terms.  For a supported vector in
\(\operatorname{Dom}a(D)\), pairing its multiplier with a form-domain
test is an \(L^2\) pairing.  The representation theorem for closed
semibounded forms therefore puts that vector in
\(\operatorname{Dom}A_L\), with
\[
\begin{split}
 A_Lv(x)={}&1_{I_L}(x)a(D)v(x)
 -1_{I_L}(x)\sum_{n\geq2}\frac{\Lambda(n)}{\sqrt n}
                    \bigl(v(x+\log n)+v(x-\log n)\bigr)\\
 &+1_{I_L}(x)\int_{\mathbb R}
       (e^{(x-y)/2}+e^{-(x-y)/2})v(y)\,dy .
\end{split}                                                    \tag{SB.5}
\]
For supported \(v\), the sum in (SB.5) has only \(n\leq\lambda^2\).
In particular no unproved endpoint boundary condition is being imposed.

By AC-S04 and elementary Gaussian tail integration,
\[
 \|r_{j,L}\|_1=O_j(\lambda^{d_j-2}e^{-\pi\lambda^2}),
 \qquad
 \operatorname{Var}(r_{j,L})
       =O_j(\lambda^{d_j}e^{-\pi\lambda^2}).                   \tag{SB.6}
\]
The variation includes the two cutoff jumps.  To check its derivative
term, integrate the AC-S04 bound for \(K^{(j+1)}\); its power is
\(d_j+2\), and integration with \(dx=du/u\) lowers it by two.
Equation (SB.4) and Plancherel now give
\[
 \|a(D)r_{j,L}\|_2
       =O_j(\lambda^{d_j}e^{-\pi\lambda^2}).                   \tag{SB.7}
\]
This controls the endpoint logarithmic singularities in \(L^2\).

Here is a direct bound for the infinite prime tail.  For
\(-L\leq x\leq L\), write \(u=e^x\in[\lambda^{-1},\lambda]\).
The positive shift can contribute only when \(n>\lambda/u\).
Since \(\Lambda(n)\leq\log n\), its absolute value is at most
\[
 C_j u^{d_j}\!
 \sum_{n>\lambda/u}\log(2n)n^{d_j-1/2}
                              e^{-\pi u^2n^2}.                \tag{SB.8}
\]
For sufficiently large \(\lambda\), the summand as a function of
the real variable \(n\) is decreasing on \(n\geq N:=\lambda/u\).
Its sum is bounded by its value at \(N\) plus its integral from
\(N\) to infinity.  After multiplication by \(u^{d_j}\), these
two quantities are at most
\[
 C_j u^{1/2}\lambda^{d_j-1/2}
                  \log(2\lambda/u)e^{-\pi\lambda^2},
 \qquad
 C_j u^{-1/2}\lambda^{d_j-3/2}
           (1+\log\lambda+|\log u|)e^{-\pi\lambda^2}.
\]
Both are \(O_j(\lambda^{d_j}\log\lambda\,e^{-\pi\lambda^2})\),
uniformly in the indicated interval.  Reflection gives the same
bound for the negative shift, since
\(|K^{(j)}(-y)|=|K^{(j)}(y)|\).
The prime series is therefore absolutely convergent on the window,
and its \(L^2\) norm is
\[
 O_j(\lambda^{d_j}(\log\lambda)^{3/2}e^{-\pi\lambda^2}).
                                                               \tag{SB.9}
\]

For the polar tail, Gaussian integration gives
\[
 \int e^{|y|/2}|r_{j,L}(y)|\,dy
     =O_j(\lambda^{d_j-3/2}e^{-\pi\lambda^2}).
\]
Multiplication by \(e^{|x|/2}\leq\lambda^{1/2}\) on the window
shows that its \(L^2\) norm is
\(O_j(\lambda^{d_j-1}\sqrt{\log\lambda}\,e^{-\pi\lambda^2})\).

Finally apply (SB.2) to \(K^{(j)}=f_{j,L}+r_{j,L}\).
The multiplier term is linear on its whole-line operator domain;
the prime series just shown is absolutely convergent on the window;
and the polar integrals converge absolutely.  Thus (SB.5) gives
\[
 A_Lf_{j,L}=-1_{I_L}\bigl(\mathcal W*r_{j,L}\bigr)
\]
as an \(L^2\) identity.  On this nonsmooth tail the notation means
the three convergent component expressions just estimated; it does
not invoke evaluation of a distribution on a discontinuous test.
The three estimates above prove (SB.3).
\end{proof}

\begin{theorem}[AC-S24: unconditional near-zero saturation and resolvent limit]
\label{AC-S24}
Define
\[
 \epsilon_\lambda=e^{-\pi\lambda^2/2},\qquad
 P_\lambda^{\rm abs}=1_{[-\epsilon_\lambda,\epsilon_\lambda]}(A_L)
\]
on \(L^2(I_L)\), and extend this finite-rank projection by zero
on the complement.  Unconditionally,
\[
 P_\lambda^{\rm abs}\longrightarrow I
       \quad\text{strongly on }L^2(\mathbb R),\qquad
 \operatorname{rank}P_\lambda^{\rm abs}\longrightarrow\infty .
                                                               \tag{SB.10}
\]
For every fixed \(j\geq0\) and \(b\geq0\),
\[
 \|P_\lambda^{\rm abs}K^{(j)}-K^{(j)}\|_{L^1(e^{b|x|}dx)}
 =O_{j,b}\!\left(
 \lambda^{b+9/2+2j}(\log\lambda)^2e^{-\pi\lambda^2/2}\right).
                                                               \tag{SB.11}
\]
In particular their Fourier transforms converge uniformly on every
fixed closed horizontal strip to \((iz)^j\Xi(z)\).
The self-adjoint operators
\[
 B_L=A_L\oplus0\quad\text{on }
        L^2(I_L)\oplus L^2(I_L^c)
\]
converge to the zero operator in strong resolvent sense.

Every family of normalized eigenvectors whose eigenvalues stay a
fixed positive distance from zero converges weakly to zero, after
extension by zero outside its window.  In particular, the following
implication is unconditional: if RH were false, every choice of
normalized bottom eigenvectors would converge weakly to zero.
This is an implication, not an assertion that such a zero exists.
\end{theorem}
\begin{proof}
The spectral theorem and (SB.3) imply
\[
 \|(I-P_\lambda^{\rm abs})f_{j,L}\|_2
 \leq\epsilon_\lambda^{-1}\|A_Lf_{j,L}\|_2
 =O_j(\lambda^{d_j}(\log\lambda)^{3/2}
                                      e^{-\pi\lambda^2/2}).
\]
Here \(P_\lambda^{\rm abs}K^{(j)}
      =P_\lambda^{\rm abs}f_{j,L}\).
Cauchy--Schwarz on the window contributes at most
\(\lambda^b\sqrt{2L}\) to the weighted \(L^1\) norm.
Outside it, AC-S04 gives a smaller Gaussian tail.  This proves
(SB.11) and ordinary \(L^2\) convergence on each derivative.

The derivatives have dense span, by the unconditional part of the
argument in AC-S18: \(|\Xi(t)|^2dt\) has exponential moments of all
rates below \(\pi/2\), and a vector orthogonal to all polynomial
multiples of \(\Xi\) gives a holomorphic Fourier transform with all
derivatives zero at the origin.  Fourier uniqueness makes that
vector zero.  Projections have norm at most one, so convergence on
the dense span proves strong convergence on the whole space.
Their ranks tend to infinity: for any \(r\), strong convergence on
\(r\) orthonormal vectors makes their images linearly independent
for all sufficiently large \(\lambda\).
Compact resolvent of \(A_L\) gives finite rank at every fixed
\(\lambda\).  No sign condition on \(A_L\) was used.

Every \(K^{(j)}\) belongs to \(\operatorname{Dom}B_L\), by AC-S23,
and \(\|B_LK^{(j)}\|_2\to0\).  For their finite linear combinations
and \(z\notin\mathbb R\), the resolvent identity gives
\[
 (B_L-z)^{-1}v+z^{-1}v
      =z^{-1}(B_L-z)^{-1}B_Lv\longrightarrow0 .
\]
The uniform bound \(\|(B_L-z)^{-1}\|\leq|\operatorname{Im}z|^{-1}\)
and density prove the asserted strong resolvent limit.

If \(u_L\) is a normalized eigenvector with
\(|a_L|\geq c>0\), it is orthogonal to
\(\operatorname{Ran}P_\lambda^{\rm abs}\) for large \(\lambda\).
For every fixed \(v\in L^2(\mathbb R)\),
\[
 |\langle u_L,v\rangle|
 \leq\|(I-P_\lambda^{\rm abs})v\|_2\longrightarrow0 .
\]
If RH were false, Weil's criterion would supply a fixed compact
smooth \(g\) with \(q(g)<0\).  For every sufficiently large window
its Rayleigh quotient is the same number
\(-c=q(g)/\|g\|_2^2<0\), hence \(a_1(L)\leq-c\).
The preceding argument applies to every bottom eigenvector.
\end{proof}

AC-S24 does not select the bottom eigenspace.  Under \(A\), the
narrower positive-energy band of AC-S18 remains available.  Without
\(A\), the symmetric band in (SB.10) still approximates every fixed
\(L^2\) vector.  Thus small operator residual and Fourier convergence
of vectors in a near-zero subspace are insufficient to identify
the bottom eigenline sought in S25 and L26.

\subsection{Primitive Dirichlet theta kernels, including cubic Dirichlet characters}

\begin{theorem}[AC-S25: the unconditional Dirichlet extension]
\label{AC-S25}
Let \(\chi\) be a primitive nonprincipal Dirichlet character of
conductor \(q\), with \(\chi(-1)=(-1)^\epsilon\).  Define
\[
 \Theta_\chi(x)=
  2e^{(1/2+\epsilon)x}
       \sum_{n\geq1}\chi(n)n^\epsilon
                     e^{-\pi n^2e^{2x}/q},\qquad
 \Lambda_\chi(s)=
 (q/\pi)^{(s+\epsilon)/2}
       \Gamma((s+\epsilon)/2)L(s,\chi).
                                                               \tag{SB.12}
\]
For \(A_{\chi,\lambda}\), use the closed semibounded operator
constructed in the proof of AC-T06.  Its compact resolvent also
follows there without assuming \(A\).
Put \(d_{\chi,j}=1/2+\epsilon+2j\) and \(\kappa_\chi=\pi/q\).
Then
\[
 1_{I_L}\Theta_\chi^{(j)}
             \in\operatorname{Dom}A_{\chi,\lambda},\qquad
 \|A_{\chi,\lambda}(1_{I_L}\Theta_\chi^{(j)})\|_2
   =O_{\chi,j}\!\left(
       \lambda^{d_{\chi,j}}(\log\lambda)^{3/2}
                                      e^{-\kappa_\chi\lambda^2}\right).
                                                               \tag{SB.13}
\]
The projection onto
\(|A_{\chi,\lambda}|\leq e^{-\kappa_\chi\lambda^2/2}\), extended
by zero, tends strongly to identity and has diverging rank.
For its projection of \(\Theta_\chi^{(j)}\), the analogue of
(SB.11) has bound
\[
 O_{\chi,j,b}\!\left(
 \lambda^{b+d_{\chi,j}}(\log\lambda)^2
                               e^{-\kappa_\chi\lambda^2/2}\right).
                                                               \tag{SB.14}
\]
The padded operators \(A_{\chi,\lambda}\oplus0\) tend to zero
in strong resolvent sense.  All these statements are unconditional.
They apply in particular to primitive cubic Dirichlet characters.
This Dirichlet theorem alone does not apply to
\(\chi_{269,3}\) over \(K\); its degree-two Hecke setting requires
the separate construction in AC-S27.
\end{theorem}
\begin{proof}
Twisted Poisson summation, as proved in (TB.6), applied to
\(x^\epsilon e^{-\pi x^2/q}\), gives
\[
 \Theta_\chi(x)=\omega_\chi\Theta_{\overline\chi}(-x),
 \qquad|\omega_\chi|=1 .
\]
There is no constant term because \(\chi\) is nonprincipal.
Direct Gaussian estimates at \(+\infty\), and this identity at
\(-\infty\), show
\[
 |\Theta_\chi^{(j)}(x)|
 \leq C_{\chi,j}e^{d_{\chi,j}|x|}
                        e^{-\kappa_\chi e^{2|x|}}
 \quad(|x|\text{ sufficiently large}).
                                                               \tag{SB.15}
\]
For \(\operatorname{Re}s>1\), termwise Mellin integration gives
\[
 \int_{\mathbb R}\Theta_\chi(x)e^{(s-1/2)x}\,dx
       =\Lambda_\chi(s).
\]
Both sides are entire, so the equality holds everywhere.
In particular
\(\widehat{\Theta_\chi}(t)=\Lambda_\chi(1/2-it)\).

For clarity about complex characters, let
\(\mathcal W_\chi\) denote the distribution in (T.11).
The actual operator in (T.8) acts, before compression, as
\[
 u(x)\longmapsto
 a_\chi(D)u(x)-
 \sum_{n\geq2}\frac{\Lambda(n)}{\sqrt n}
 \bigl(\chi(n)u(x+\log n)+
             \overline{\chi(n)}u(x-\log n)\bigr).
                                                               \tag{SB.16}
\]
This is \(\mathcal W_{\chi,y}(u(x+y))\), rather than a
convolution with the unreversed character distribution.
The extension argument in AC-S22 applies to the Dirichlet
explicit formula, since its zero count is
\(O_\chi(T\log(2+T))\).
For \(u=\Theta_\chi^{(j)}\), the zero-side term at \(\rho\) is
\[
 e^{-(\rho-1/2)x}(-(\rho-1/2))^j\Lambda_\chi(\rho)=0 .
\]
Thus (SB.16) annihilates every derivative without using GRH.

The proof of AC-S23 now applies with \(d_j\) replaced by
\(d_{\chi,j}\) and \(\pi\) by \(\kappa_\chi\).
The multiplier \(a_\chi\) still has logarithmic growth.
The prime terms have the same absolute bounds because
\(|\chi(n)|\leq1\).  There are no polar terms.
The variation and endpoint argument is unchanged, and the
Gaussian tail bounds (SB.15) prove (SB.13).

The derivative span is dense.  Bounded partial sums of a
nonprincipal Dirichlet character give
\(L(1/2+it,\chi)=O_\chi(1+|t|)\) by partial summation.
Stirling's formula therefore gives exponential moments, with
every rate below \(\pi/2\), for
\(|\Lambda_\chi(1/2-it)|^2dt\).
Analytic Fourier uniqueness proves density exactly as in AC-S24.
The spectral and weighted-norm arguments of that theorem prove
all the remaining assertions, with the displayed constants and
exponents.  Reflection or evenness of the complex-character
operator is not assumed.
\end{proof}

\subsection{A concrete obstruction to ordinary cone positivity}

\begin{proposition}[AC-S26: the Weil heat semigroup is not positivity preserving]
\label{AC-S26}
If \(\lambda>\sqrt{3/2}\), the semigroup \(e^{-sA_L}\),
\(s>0\), is not positivity preserving on the ordinary cone
of nonnegative functions in \(L^2(I_L)\).
If \(\lambda>3/2\), its restriction to the even subspace is
also not positivity preserving on the ordinary even cone.
In each case the failure occurs for all sufficiently small
positive \(s\).
These statements are unconditional.  They do not disprove
simplicity or evenness of the lowest eigenvector.
\end{proposition}
\begin{proof}
Away from the diagonal and from the prime translation atoms,
(SB.1) and S21's \texttt{bombieriexplicit2bis} give the
off-diagonal kernel
\[
 k(t)=2\cosh(t/2)-
       \frac{e^{|t|/2}}{e^{|t|}-e^{-|t|}},\qquad t\ne0 .
                                                               \tag{SB.17}
\]
The first summand is exactly the two polar terms; the second is
the off-diagonal archimedean term.  The subtraction at the origin
in the principal value makes no contribution for disjoint supports.
For \(r=e^{|t|}>1\),
\[
 k(t)=\frac{r^3-r-1}{\sqrt r\,(r^2-1)} .
\]
In particular, with \(a=\log(3/2)<\log2\),
\[
 k(a)=\frac7{10\sqrt{3/2}}>0 .
\]
If \(L>a/2\), choose nonzero nonnegative smooth bumps \(f,g\)
in small neighborhoods of \(-a/2,a/2\), respectively.
The neighborhoods can be chosen inside \(I_L\), disjoint, with
all separation distances in \((0,\log2)\) and in a small
neighborhood of \(a\) on which \(k\) is positive.
No prime translation atom has such a distance.  Therefore
\[
 \langle f,A_Lg\rangle
     =\iint f(x)g(y)k(x-y)\,dx\,dy>0 .
\]
Both bumps belong to the operator domain, by the same multiplier
argument as in AC-S23.  Since \(\langle f,g\rangle=0\),
strong differentiability of the semigroup on
\(\operatorname{Dom}A_L\) at \(s=0\) gives
\[
 \langle f,e^{-sA_L}g\rangle
       =-s\langle f,A_Lg\rangle+o(s)<0
\]
for sufficiently small positive \(s\).
A positivity-preserving operator cannot have a negative pairing
between two nonnegative functions.  This proves the first claim.

For the even claim, assume \(L>a\), take \(f\) to be an even bump
near zero, and \(g\) the sum of two reflected nonnegative bumps
near \(a\) and \(-a\).  Every relevant separation is again near
\(\pm a\), so the same strict positive form pairing and negative
semigroup pairing result.  Reflection commutes with the zeta
operator, so this is also a counterexample inside its even sector.
\end{proof}

The usual Perron--Frobenius argument based on preservation of the
ordinary nonnegative cone cannot therefore establish the desired
growing-window ground-state property.  A different cone,
oscillation theorem, or quantitative comparison may still succeed.
The remaining bottom-eigenline gate in AC-S07 is left explicit;
neither the near-zero quasimodes nor this cone obstruction settles
that gate.

\subsection{The degree-two Hecke construction}

The primitive Hecke character of a cubic Kummer extension of the
imaginary quadratic field \(K\) is a different input from a cubic
Dirichlet character over \(\mathbb Q\).  The following construction
treats the former directly.  Its classical background is analytic
continuation and the functional equation for a primitive finite-order
Hecke character.  The exact completion used here is also printed in
Brattstr\"om--Lichtenbaum, \emph{Jacobi-sum Hecke characters of imaginary
quadratic fields}, \emph{Compositio Math.} 53 (1984), Theorem 4.3,
p.~294.  The new consequences below are proved from that classical
input; no RH for these Hecke functions is assumed.

\begin{theorem}[AC-S27: unconditional Hecke theta quasimodes in degree two]
\label{AC-S27}
Let \(K\) be imaginary quadratic and let \(\chi\) be a primitive,
nonprincipal, finite-order Hecke character, of exact conductor
\(\mathfrak f\).  Put
\[
 Q_\chi=|D_K|N\mathfrak f,\qquad c_\chi=\frac{2\pi}{\sqrt{Q_\chi}},
 \qquad
 a_\chi(n)=\sum_{N\mathfrak a=n}\chi(\mathfrak a),
\]
where the character is zero on ideals not coprime to its conductor.
Define
\[
 \Lambda_\chi(s)=c_\chi^{-s}\Gamma(s)L_K(s,\chi),\qquad
 \Theta_\chi(x)=e^{x/2}
             \sum_{n\geq1}a_\chi(n)e^{-c_\chi n e^x}.
                                                               \tag{SB.18}
\]
The completion omits the harmless constant factor \(2\) sometimes
included in \(\Gamma_{\mathbb C}\).

There is an unconditional closed lower semibounded Weil operator
\(A^K_{\chi,\lambda}\) on \(L^2(I_L)\), with compact resolvent,
whose archimedean symbol and prime-ideal translations are
\[
 a^K_\chi(t)=-2\log c_\chi+
                  2\operatorname{Re}\psi(1/2+it),               \tag{SB.19}
\]
\[
 -\sum_{\mathfrak p,k\geq1}
 \frac{\log N\mathfrak p}{(N\mathfrak p)^{k/2}}
 \bigl(\chi(\mathfrak p)^kT_{-k\log N\mathfrak p}
       +\overline{\chi(\mathfrak p)^k}T_{k\log N\mathfrak p}\bigr),
                                                               \tag{SB.20}
\]
where only \((N\mathfrak p)^k\leq\lambda^2\) survive compression.
There are no polar terms.
For every fixed \(j\geq0\),
\[
 1_{I_L}\Theta_\chi^{(j)}
        \in\operatorname{Dom}A^K_{\chi,\lambda},\qquad
 \|A^K_{\chi,\lambda}(1_{I_L}\Theta_\chi^{(j)})\|_2
 =O_{\chi,j}\!\left(
       \lambda^{1/2+j}(\log\lambda)^{3/2}e^{-c_\chi\lambda}\right).
                                                               \tag{SB.21}
\]
The finite-rank projections
\[
 P^K_{\chi,\lambda}
   =1_{[-e^{-c_\chi\lambda/2},\,e^{-c_\chi\lambda/2}]}
                                      (A^K_{\chi,\lambda}),
\]
extended by zero, tend strongly to identity and have diverging rank.
Their projected theta derivatives satisfy
\[
 \|P^K_{\chi,\lambda}\Theta_\chi^{(j)}
                   -\Theta_\chi^{(j)}\|_{L^1(e^{b|x|}dx)}
 =O_{\chi,j,b}\!\left(
 \lambda^{b+1/2+j}(\log\lambda)^2e^{-c_\chi\lambda/2}\right)
 \quad(b\geq0).                                                \tag{SB.22}
\]
The operators \(A^K_{\chi,\lambda}\oplus0\) tend to zero in
strong resolvent sense.  This applies to the primitive Hecke
character underlying \(\chi_{269,3}\), with its exact conductor.
It neither proves its RH nor identifies a bottom eigenline.
\end{theorem}
\begin{proof}
At each rational prime, the number of quadratic-field ideals of
norm \(p^m\) is at most \(m+1\).  Unique ideal factorization therefore
gives \(|a_\chi(n)|\leq d(n)\), where \(d(n)\) is the divisor
function.  The series in (SB.18) and all its derivatives converge
locally uniformly.  For \(x\) large and \(u=e^x\),
\[
 |\Theta_\chi^{(j)}(x)|
       \leq C_{\chi,j}u^{1/2+j}e^{-c_\chi u}.                  \tag{SB.23}
\]
For example, differentiating produces a polynomial of degree at
most \(j\) in \(nu\); after factoring out \(e^{-c_\chi u}\),
the remaining series is bounded uniformly for \(u\geq2\).

Termwise Mellin integration on \(\operatorname{Re}s>1\) gives
\[
 \int_{\mathbb R}\Theta_\chi(x)e^{(s-1/2)x}\,dx
       =\Lambda_\chi(s).
\]
Use the classical entire continuation and functional equation
\[
 \Lambda_\chi(s)=\omega_\chi\Lambda_{\overline\chi}(1-s),
 \qquad|\omega_\chi|=1.
                                                               \tag{SB.24}
\]
Mellin inversion on a line \(\operatorname{Re}s=\sigma>1\),
substitution of (SB.24), and the change of variable \(s\mapsto1-s\)
give
\(\Theta_\chi(x)=\omega_\chi\Theta_{\overline\chi}(-x)\).
Here the contour can be moved back from \(1-\sigma\) to \(\sigma\):
the completed function is entire; the Dirichlet series,
functional equation, and polynomial convexity give polynomial
growth of \(L_K(s,\chi)\) on fixed vertical strips; Stirling's
formula for \(\Gamma(s)\) makes the horizontal integrals tend
to zero exponentially.  Thus no residues are omitted.
Equation (SB.23) and this identity give at both ends
\[
 |\Theta_\chi^{(j)}(x)|
 \leq C_{\chi,j}e^{(1/2+j)|x|}
                              e^{-c_\chi e^{|x|}}.
                                                               \tag{SB.25}
\]
In particular the bilateral Mellin integral is entire, and the
identity above holds for all \(s\).

The general Hecke explicit formula, C99, Appendix II, Theorem 6,
has symbol twice the real logarithmic derivative of the gamma
and conductor factor, which is exactly (SB.19).
Its remaining terms are (SB.20).  The symbol is real, bounded
below, and \(2\log|t|+O_\chi(1)\) at infinity.  Restricting its
closed weighted Fourier form to \(I_L\), then adding the finite
bounded Hermitian translation sum, gives a closed semibounded
form.  The Fourier-tail compactness argument of AC-T06 proves
compact resolvent; its smooth-core argument applies unchanged.
This establishes the asserted operator independently of any
zero-location hypothesis.

The extension of the explicit formula in AC-S22 applies to
(SB.25), using the usual \(O_\chi(T\log(2+T))\) Hecke zero count.
As in (SB.16), apply the character distribution to
\(\Theta_\chi^{(j)}(x+y)\).
Every zero-side term is
\[
 e^{-(\rho-1/2)x}(-(\rho-1/2))^j\Lambda_\chi(\rho)=0.
 \]
Thus the full distribution annihilates each theta derivative
with precisely the orientation of the translations in (SB.20).

For completeness, the change from \(e^{-\pi e^{2|x|}}\) to
\(e^{-c_\chi e^{|x|}}\) changes the tail estimates as follows.
Put \(d=1/2+j\) and
\(r=(1-1_{I_L})\Theta_\chi^{(j)}\).
Equation (SB.25) gives
\[
 \|r\|_1=O_{\chi,j}(\lambda^{d-1}e^{-c_\chi\lambda}),
 \qquad
 \operatorname{Var}(r)=O_{\chi,j}(\lambda^de^{-c_\chi\lambda}).
\]
Consequently the multiplier tail has \(L^2\) norm
\(O_{\chi,j}(\lambda^de^{-c_\chi\lambda})\), by (SB.4).
The sharp cutoff belongs to the operator domain by the same
form-representation argument.

Group the prime-ideal coefficients according to
\(n=(N\mathfrak p)^k\).  Their absolute sum at a fixed \(n\)
is at most \(2\log n\): a split rational prime contributes
at most two terms \(\log p\), an inert prime at most one
term \(2\log p\), and a ramified prime at most one term
\(\log p\).  For \(u=e^x\in[\lambda^{-1},\lambda]\), the
positive-shift tail is therefore bounded by
\[
 C_{\chi,j}u^d\sum_{n>\lambda/u}
               \log(2n)n^{d-1/2}e^{-c_\chi un}.
\]
Its decreasing-summand value at \(N=\lambda/u\) and its
integral from \(N\) to infinity are bounded respectively by
\[
 C_{\chi,j}u^{1/2}\lambda^{d-1/2}
                  \log(2\lambda/u)e^{-c_\chi\lambda},
 \quad
 C_{\chi,j}u^{-1/2}\lambda^{d-1/2}
       (1+\log\lambda+|\log u|)e^{-c_\chi\lambda}.
\]
Both are \(O_{\chi,j}(\lambda^d\log\lambda
e^{-c_\chi\lambda})\).
The negative-shift tail has the same bound by (SB.25).
Multiplication by \(\sqrt{2L}\) gives (SB.21).
No polar tail occurs.  The cancellation with the global
null identity is interpreted through these convergent
component expressions, as in AC-S23.

Finally
\(\widehat{\Theta_\chi}(t)=\Lambda_\chi(1/2-it)\).
Polynomial growth of \(L_K\) and Stirling's formula show that
\(|\widehat{\Theta_\chi}(t)|^2dt\) has exponential moments
of every rate less than \(\pi\).
The nonzero entire completion has only a discrete real zero
set, so analytic Fourier uniqueness again makes the theta
derivatives total in \(L^2(\mathbb R)\).
The spectral estimate, weighted Cauchy--Schwarz, and resolvent
identity in AC-S24 now prove (SB.22), strong convergence of
the projections, rank divergence, and the padded resolvent
limit, with \(\epsilon_\lambda=e^{-c_\chi\lambda/2}\).
\end{proof}

\subsection{A proved small-window ground eigenline}

\begin{theorem}[AC-S28: small-window simplicity, parity, and the exact cone threshold]
\label{AC-S28}
Let \(r_0>1\) be the real root of \(r^3-r-1=0\).
For
\[
                 1<\lambda\leq\sqrt{r_0},
\]
the actual zeta Weil operator \(A_L\), \(L=\log\lambda\), has a
one-dimensional lowest eigenspace with a nonnegative even generator.
Its heat semigroup preserves the ordinary nonnegative cone.
For \(\lambda>\sqrt{r_0}\), that semigroup fails to preserve this
cone for all sufficiently small positive times.
All these assertions are unconditional.  The theorem does not
establish simplicity or evenness in the growing-window regime.
\end{theorem}
\begin{proof}
Put \(I=(-L,L)\), extend vectors by zero outside \(I\), and set
\[
 \mu(t)=\frac{e^{-|t|/2}}{1-e^{-2|t|}},\qquad
 p(t)=2\cosh(t/2),\qquad J(t)=\mu(t)-p(t),\quad t\ne0.
\]
The integral formula for the digamma difference gives
\[
 a(\tau)=a(0)+2\int_0^\infty
                 \mu(t)(1-\cos(\tau t))\,dt .
\]
Plancherel and Tonelli therefore identify the archimedean form with
\[
 a(0)\|u\|_2^2+
 \frac12\iint_{\mathbb R^2}
                \mu(x-y)|u(x)-u(y)|^2\,dx\,dy.
                                                               \tag{SB.26}
\]
This identity holds on its full closed form domain: the nonnegative
integral is exactly the Fourier integral for \(a(\tau)-a(0)\),
including the value \(+\infty\) before restriction to the domain.

When \(\lambda^2\leq r_0<2\), the interval has length
\(2L<\log2\), so every compressed prime translation is zero.
Using
\[
 \begin{aligned}
 &\iint_{I^2}p(x-y)\overline{u(x)}u(y)\,dx\,dy\\
 &\quad=\int_I\left(\int_Ip(x-y)\,dy\right)|u(x)|^2\,dx
       -\frac12\iint_{I^2}p(x-y)|u(x)-u(y)|^2\,dx\,dy,
 \end{aligned}
\]
where symmetry makes the left side real, the full form becomes
\[
 \begin{aligned}
 q_L(u)&=\int_I V(x)|u(x)|^2\,dx+
       \frac12\iint_{I^2}J(x-y)|u(x)-u(y)|^2\,dx\,dy,
 \\
 V(x)&=a(0)+\int_{\mathbb R\setminus I}\mu(x-y)\,dy
                         +\int_Ip(x-y)\,dy .
 \end{aligned}                                                \tag{SB.27}
\]
Here \(V\geq a(0)\); it is finite in the interior and has only
logarithmic endpoint divergence.  Subtracting the bounded kernel
\(p\) does not alter the domain in (SB.26), so (SB.27) is an
identity on the actual full closed form domain.

The algebra in (SB.17) shows
\[
 J(t)>0\qquad(0<|t|<\log r_0).
\]
Consequently \(J(x-y)>0\) for almost every pair of distinct points
in \(I\), also at the endpoint parameter \(2L=\log r_0\).
The elementary inequality
\(\bigl||u(x)|-|u(y)|\bigr|\leq|u(x)-u(y)|\) proves that
\(|u|\) belongs to the form domain and \(q_L(|u|)\leq q_L(u)\).
For a real vector with positive and negative parts both nonzero,
the inequality is strict: on their Cartesian product the decrease
of the squared difference is \(4|u(x)u(y)|\), and \(J>0\)
almost everywhere.  All integrals are finite on the form domain.

Compact resolvent and lower semiboundedness give a lowest eigenvalue
and an eigenvector.  Reality of (SB.27) permits real eigenvectors.
The variational principle and the strict inequality just proved
show that every nonzero real lowest eigenvector has one sign.
If this eigenspace had real dimension at least two, take two real
orthonormal eigenvectors and change their signs to make both
nonnegative.  Orthogonality makes their positive supports disjoint;
their difference is a lowest eigenvector with both signs, a
contradiction.  Thus the complex eigenspace is one-dimensional.
Reflection commutes with the form.  It preserves the nonnegative
unit generator, so that generator is even.

For completeness, cone preservation follows directly from the same
modulus inequality.  If \(c>-a(0)\), the strictly convex resolvent
functional
\(q_L(u)+c\|u\|_2^2-2\operatorname{Re}\langle f,u\rangle\),
for \(f\geq0\), does not increase upon replacing \(u\) by \(|u|\).
Its unique minimizer is therefore nonnegative.  The resolvent
\((A_L+c)^{-1}\) is positive.  Applying this to all larger shifts
and using the scalar spectral-calculus Euler limit for the heat
semigroup proves cone preservation.

Finally, if \(2L>\log r_0\), choose
\(a\in(\log r_0,\min(2L,\log2))\).
Then \(k(a)=-J(a)>0\), and no prime atom occurs at a nearby
distance.  The disjoint bumps near \(-a/2\) and \(a/2\) in
AC-S26 give a strictly negative heat-semigroup pairing for all
sufficiently small positive times.  This proves the asserted
failure and the exact threshold.  It does not imply a failure of
simplicity or parity beyond the threshold.
\end{proof}

\begin{remark}[The numerical limit on the proved parameter range]
For each fixed \(1<\lambda\leq\sqrt{r_0}\), AC-S28 removes the
simplicity and evenness hypotheses in AC-S06.  Thus, after phase
and \(L^2\) normalization, the entire sequence of lowest Galerkin
eigenvectors converges in form norm and its Fourier transforms
converge locally uniformly.  The finite matrices are even-simple
for all sufficiently large dimensions, and AC-S06's normalized
determinant conclusion applies.  This consequence is unconditional:
the Galerkin minimax and compactness argument only requires lower
semiboundedness and the form core; adding a fixed scalar to the
form makes its norm positive without using \(A\).
It is a fixed-window conclusion and asserts neither convergence
of the boundary normalization nor the large-window Xi limit.
\end{remark}
